%%%%%%%%%%%%%%%%%%%%%%%%%%%%%%%%%%%%%%%%%
% Developer CV
% LaTeX Class
% Version 2.0 (12/10/23)
%
% This class originates from:
% http://www.LaTeXTemplates.com
%
% Authors:
% Omar Roldan
% Based on a template by  Jan Vorisek (jan@vorisek.me)
% Based on a template by Jan Küster (info@jankuester.com)
% Modified for LaTeX Templates by Vel (vel@LaTeXTemplates.com)
%
% License:
% The MIT License (see included LICENSE file)
%
%%%%%%%%%%%%%%%%%%%%%%%%%%%%%%%%%%%%%%%%%

%----------------------------------------------------------------------------------------
%	PACKAGES AND OTHER DOCUMENT CONFIGURATIONS
%----------------------------------------------------------------------------------------

\documentclass[9pt]{developercv} % Default font size, values from 8-12pt are recommended
\usepackage{multicol}
\setlength{\columnsep}{0mm}
%----------------------------------------------------------------------------------------
\usepackage{lipsum}  


\begin{document}

%----------------------------------------------------------------------------------------
%	TITLE AND CONTACT INFORMATION
%----------------------------------------------------------------------------------------

\begin{minipage}[t]{0.5\textwidth} 
	\vspace{-\baselineskip} % Required for vertically aligning minipages
	
	{ \fontsize{16}{20} \textcolor{black}{\textbf{\MakeUppercase{Hrishikesh Belagali}}}} % First name
	
	\vspace{6pt}
	
	{\Large Undergraduate Student} % Career or current job title
\end{minipage}
\hfill
\begin{minipage}[t]{0.2\textwidth} % 20% of the page width for the first row of icons
	\vspace{-\baselineskip} % Required for vertically aligning minipages
	
	% The first parameter is the FontAwesome icon name, the second is the box size and the third is the text
	% \icon{Globe}{11}{\href{http://www.google.com}{portafolio.com}}\\ 
    % \icon{Phone}{11}{11 1111 1111}\\
    \icon{MapMarker}{11}{East Lansing}\\
	\icon{Envelope}{11}{\href{mailto:belagal1@msu.edu}{belagal1@.msu.edu}}\\	
\end{minipage}
\begin{minipage}[t]{0.27\textwidth} % 27% of the page width for the second row of icons
	\vspace{-\baselineskip} % Required for vertically aligning minipages
    \icon{Github}{11}{\href{https://github.com/lonelyneutrin0}{github.com/lonelyneutrin0}}\\
    \icon{LinkedinSquare}{11}{\href{https://www.linkedin.com/in/hkbel}{/in/hkbel}}\\    
    
\end{minipage}


%----------------------------------------------------------------------------------------
%	INTRODUCTION, SKILLS AND TECHNOLOGIES
%----------------------------------------------------------------------------------------

\begin{minipage}[t]{0.46\textwidth}
\cvsect{Summary}
\vspace{-6pt}

Undergraduate computer science major with research interests in quantum algorithms, molecular dynamics, reservoir computing, quantum dynamics, and high performance computing.\\
\end{minipage}
\hfill % Whitespace between
\begin{minipage}[t]{0.465\textwidth}
    \cvsect{Skills}
\\\begin{minipage}[t]{0.2\textwidth}
        \textbf{Languages:}
    \end{minipage}
    \hfill
    \begin{minipage}[t]{0.73\textwidth}
      Python, C, C++, \LaTeX
    \end{minipage}
    \vspace{4mm}
    
    \begin{minipage}[t]{0.2\textwidth}
        \textbf{Packages:}
    \end{minipage}
    \hfill
    \begin{minipage}[t]{0.73\textwidth}
    NumPy, PyTorch, scipy, Qiskit, Cirq, PyEMMA, hyperopt, VMD-Python
    \end{minipage}
	\vspace{4mm}

	\begin{minipage}[t]{0.2\textwidth}
        \textbf{Softwares:}
    \end{minipage}
    \hfill
    \begin{minipage}[t]{0.73\textwidth}
    LAMMPS, OVITO, NAMD, VMD
    \end{minipage}
    \vspace{1mm}
\end{minipage}

%----------------------------------------------------------------------------------------
%	EDUCATION
%----------------------------------------------------------------------------------------
\vspace{-10 pt}
\cvsect{Education}
\begin{entrylist}
    \entry
		{\hspace{10pt}}
		{Computer Science (B.S)}
		{Michigan State University Honors College}
		{Expected Graduation: 2027 \\ 
		Current GPA: 4.0}
\end{entrylist}
\vspace{-10pt}
\cvsect{EXPERIENCE}
\begin{entrylist}
	\entry
	{2025--}
	{Honors Research Scholar}
	{MSU Honors College}
	{\vspace{-10pt}
	\begin{itemize}[noitemsep,topsep=0pt,parsep=0pt,partopsep=0pt, leftmargin=7pt]
	\item Conducted research on generalized quantum simulators under the guidance of Professor Ryan LaRose.
	\item Investigated computational group theory methods for subgroup membership in unitary subgroups.
	\end{itemize}
	}
	\entry 
		{2024--}
		{Undergraduate Research Assistant}
		{Vermaas Lab, Plant Research Laboratory}
		{\vspace{-10pt}
		\begin{itemize}[noitemsep,topsep=0pt,parsep=0pt,partopsep=0pt, leftmargin=7pt]
			\item Investigated the interaction of the chloroplast targeting protein (At4g16060) with thylakoid membranes using molecular dynamics.
			\item Analyzed the effect of PDLP-5 (Plasmodesmata Located Protein) on water flow rates through PIP1;4 aquaporin tetramer using grid-steered molecular dynamics.
			\item Simulated the effects of small molecules like butanol on \textit{Clostridium acetobutylicum} membrane dynamics. 
		\end{itemize}
		}
\end{entrylist}
\cvsect{CONFERENCES}
\begin{entrylist}
	\entry 
		{2025 \\ \vspace{4pt}\color{gray} \textit{PRL, MSU}}
		{Aquaporin Water Transit Modulation via Protein Capping}
		{\href{https://acs.digitellinc.com/p/s/aquaporin-water-transit-modulation-via-protein-capping-638984}{acs.digitellinc.com}}
		{Abstract presented at ACS Fall 2025 Conference}
	\entry 
		{2025 \\ \vspace{4pt}\color{gray} \textit{PRL, MSU}}
		{Evaluating the Interaction of Chloroplast-targeting Protein (At4g16060) with
		Thylakoid Membrane}
		{}
		{Abstract presented at First Great Lakes Plant Science Conference}
\end{entrylist}
%----------------------------------------------------------------------------------------
%	Projects
%----------------------------------------------------------------------------------------
\cvsect{Projects}
\begin{entrylist}
	\entry
		{Python \\ Qiskit}
		{Discretized Quantum Adiabatic Evolution}
		{\href{https://github.com/lonelyneutrin0/Adiabatic-Quantum-Computation}{github.com}}
		{Implemented a quantum adiabatic algorithm to solve QUBO type problems using the Trotter-Suzuki approximation and quantum Ising model formulation techniques}
	\entry 
		{Python \\ Qiskit}
		{Quantum Arithmetic Circuits}
		{\href{https://github.com/lonelyneutrin0/Quantum-Arithmetic-Circuits}{github.com}}	
		{Built quantum circuits to perform addition, subtraction and multiplication using the Quantum Fourier Transform.}	
	
	\entry
		{Python \\ NumPy}
		{Monte Carlo Integration}
		{\href{https://github.com/lonelyneutrin0/MNI}{github.com}}
		{Monte Carlo integration with importance sampling to numerically evaluate complex integrals.}
	\entry
		{Python \\ NumPy}
		{Langevin Monte Carlo}
		{\href{https://github.com/lonelyneutrin0/Metropolis-Langevin-Algorithm}{github.com}}
		{Metropolis-adjusted Langevin algorithm for sampling from intractable probability distributions.}
	\entry
		{Python \\ OVITO \\ pyMOL }
		{Genetic Annealing to Determine Protein Structures}
		{\href{https://github.com/lonelyneutrin0/Protein-Structure-Prediction}{github.com}}
		{Monte Carlo Genetic Algorithm to determine protein structures through the optimization of Irbäck's off-lattice model energy equation (RMSD < 3.0).}	
	% \entry
	% 	{Python \\ NumPy}
	% 	{foldopt}
	% 	{\href{https://github.com/lonelyneutrin0/foldopt}{github.com}}
	% 	{Python package for model-agnostic protein structure optimization.}
	\entry
		{Python}
		{Quadratic Unconstrained Binary Optimization}
		{\href{https://github.com/lonelyneutrin0/qubo}{github.com}}
		{Stochastic tunneling-enhanced simulated annealing for solving QUBO-type problems.}
	\entry
		{Python \\ NumPy}
		{aqua-blue}
		{\href{https://github.com/jwjeffr/aqua-blue}{github.com}/\href{https://pypi.org/project/aqua-blue/}{pypi.org}}
		{Secondary maintainer of aqua-blue, a package building echo state networks.
		\begin{itemize}[noitemsep,topsep=0pt,parsep=0pt,partopsep=0pt, leftmargin=7pt]
		\item Implemented reservoir state time-lag, sparsity and spectral radius to improve prediction accuracy.
		\item Introduced timezone-aware datetime support for seamless interfacing with data pipelines.
		\item Demonstrated prediction capability by predicting quantum time evolution of Gaussian wave packets. 
		\end{itemize}
		}
	\entry
		{Python \\ NumPy \\ hyperopt}
		{aqua-blue-hyperopt}
		{\href{https://github.com/lonelyneutrin0/aqua-blue-hyperopt}{github.com}/\href{https://pypi.org/project/aqua-blue-hyperopt/}{pypi.org}}
		{Primary maintainer of aqua-blue-hyperopt, a package implementing hyperparameter optimization functionality for aqua-blue models.
		\begin{itemize}[noitemsep,topsep=0pt,parsep=0pt,partopsep=0pt, leftmargin=7pt]
		\item Created modular hyperparameter optimization functionality for aqua-blue reservoir computers.
		\item Implemented custom loss function capabilities for flexible model evaluation.
        \end{itemize}
		}
	\entry
		{Python \\ PyTorch \\ pytest}
		{neurop}
		{\href{https://github.com/lonelyneutrin0/neurop/}{github.com}/\href{https://pypi.org/project/neurop/}{pypi.org}}
		{Primary maintainer of neurop, a package implementing neural operators.
		\begin{itemize}[noitemsep,topsep=0pt,parsep=0pt,partopsep=0pt, leftmargin=7pt]
		\item Developed Fourier Neural Operator (FNO) architecture for solving PDEs.
		\item Implemented the Fractional Fourier Transform for use in Complex Neural Operators (CoNO).
		\item Built CI/CD pipeline with GitHub Actions for automated testing and deployment.
		\end{itemize}
		}
\end{entrylist}
\cvsect{Open Source Contributions}
\begin{entrylist}
	\entry
		{QuTiP \\ Python}
		{Quantum Toolbox in Python (QuTiP)}
		{\href{https://github.com/qutip/qutip}{github.com}}
		{Contributed to QuTiP, an open-source software for simulating quantum systems.
		\begin{itemize}[noitemsep,topsep=0pt,parsep=0pt,partopsep=0pt, leftmargin=7pt]
		\item Implemented functionality to return eigenkets as an operator in the Qobj class.
		\end{itemize}
		}
\end{entrylist}
\cvsect{HONORS AND AWARDS}
\vspace{-10pt}
\begin{itemize}
	\item Wielenga Research Scholarship, Honors College (2025-2026)
	\item Dean's List (2024, 2025)
	\item Senator Joe and Mary Conroy Endowed Scholarship, College of Engineering (2024)
	\item MSU President's Scholarship (2024)
	\item Senator Joe and Mary Conroy Endowed Scholarship, College of Engineering (2024)
\end{itemize}
\end{document}
